\documentclass[11pt]{article}

\usepackage{amsmath,amssymb}
\usepackage{enumerate}
\usepackage[margin=1.2in]{geometry}
\usepackage{fancyhdr}
\usepackage{xcolor}
\usepackage{hyperref}

\pagestyle{fancy}
\setlength{\headheight}{20pt}
\setlength{\headsep}{20pt}
\fancyhead[L]{\small{CS 171, Fall 2026}}
\fancyhead[R]{\small{Prof. Sanjam Garg}}

\newcommand{\Gen}{\mathsf{Gen}}
\newcommand{\Sign}{\mathsf{Sign}}
\newcommand{\Verify}{\mathsf{Verify}}
\newcommand{\GG}{\mathbb{G}}
\newcommand{\ZZ}{\mathbb{Z}}
\newcommand{\A}{\mathcal{A}}
\newcommand{\B}{\mathcal{B}}
\newcommand{\D}{\mathcal{D}}
\newcommand{\negl}{\mathsf{negl}}
\newcommand{\duedate}{Saturday, November 14, 2026 at 8:59 PM via Gradescope}


\begin{document}
\section*{CS 171: Problem Set 7\\ {\small Due Date: \duedate}}

\section{Aggregate Signatures (15 points)}
Let $(\GG,\GG_T,q,g,e)$ be bilinear-group parameters, where $\GG$ and
$\GG_T$ have prime order $q$, $g$ generates $\GG$, and
$e:\GG\times\GG\to\GG_T$ is an efficiently computable, nondegenerate
bilinear map. Let $H:\{0,1\}^*\to\GG$ be a random oracle.
All users share these parameters and $H$.

Consider the following signature scheme $\Sigma=(\Gen,\Sign,\Verify)$:
\begin{itemize}
    \item $\Gen(1^n)$: Sample $x\leftarrow\ZZ_q$ and output
    $pk=h=g^x$ and $sk=x$.
    \item $\Sign(x,m)$: Output $\sigma=H(m)^x$.
    \item $\Verify(h,m,\sigma)$: Accept if $\sigma\in\GG$ and
    $e(g,\sigma)=e(h,H(m))$; otherwise reject.
\end{itemize}
Assume that $\Sigma$ is existentially unforgeable under chosen-message
attacks.

Alice and Bob have public keys $h_A=g^a$ and $h_B=g^b$. They sign the
same message $m$, producing $\sigma_A=H(m)^a$ and $\sigma_B=H(m)^b$.
Their \emph{aggregate signature} is the single group element
\[
    \sigma=\sigma_A\sigma_B.
\]

\begin{enumerate}[(a)]
    \item \textbf{(4 points)} Give a verification algorithm for an
    aggregate signature, given $h_A,h_B,m,\sigma$, using two evaluations
    of the bilinear map. Prove correctness.

    \item \textbf{(7 points)} Suppose Alice's and Bob's key pairs are
    generated honestly and independently. An adversary receives
    $h_A,h_B$ and Bob's secret key $b$, and may query Alice's signing
    oracle. It wins if it outputs an accepted aggregate signature on a
    message that it never queried to Alice's signing oracle.
    Prove that every PPT adversary wins with negligible probability.
    Give a reduction to the unforgeability of $\Sigma$.

    \item \textbf{(4 points)} Now suppose the adversary may choose
    $h_B\in\GG$ after receiving Alice's public key $h_A$, without
    providing the corresponding secret key. Prove that with this
    modification, the aggregate signature scheme is insecure.
\end{enumerate}


\newpage
\section{Key Exposure in a Diffie--Hellman Ratchet (15 points)}
Let $\GG$ be a cyclic group of prime order $q$ with generator $g$, in
which the DDH assumption holds. Consider five successive public values
in a Diffie--Hellman ratchet:
\[
    A_1=g^{a_1},\quad A_2=g^{a_2},\quad A_3=g^{a_3},\quad
    A_4=g^{a_4},\quad A_5=g^{a_5},
\]
where $a_1,\ldots,a_5\leftarrow\ZZ_q$ are independent and uniform.
The four successive keys are
\[
    K_{12}=g^{a_1a_2},\qquad K_{23}=g^{a_2a_3},\qquad
    K_{34}=g^{a_3a_4},\qquad K_{45}=g^{a_4a_5}.
\]
A passive adversary records all five public values and later learns
$a_3$. No other exponents or stored keys are exposed.

\begin{enumerate}[(a)]
    \item \textbf{(5 points)} Which of the four keys can the adversary
    compute? Give explicit computations for the exposed keys.

    \item \textbf{(10 points)} Let
    \[
        V=(A_1,A_2,A_3,A_4,A_5,a_3,K_{23},K_{34}).
    \]
    Prove that
    \[
        (V,K_{12},K_{45})\;\approx_c\;(V,U_1,U_2),
    \]
    where $U_1,U_2$ are independent uniform elements of $\GG$, independent
    of $V$. Give a reduction to DDH.
\end{enumerate}


\end{document}
